% ===================================================================
%                       PREÁMBULO DEL DOCUMENTO
% ===================================================================
\documentclass[12pt, a4paper]{article}

% --- PAQUETES ESENCIALES (Codificación, Idioma) ---
\usepackage[utf8]{inputenc}
\usepackage[spanish,es-lcroman]{babel}

% --- PAQUETES PARA CITAS Y BIBLIOGRAFÍA (MUY IMPORTANTE EL ORDEN) ---
\usepackage{csquotes}
\usepackage[
    backend=biber,
    style=apa,
    sortcites=true,
    maxcitenames=2,
    giveninits=true,
    uniquename=init,
    uniquelist=minyear,
    maxbibnames=99,
    date=year,
    urldate=long,
    isbn=false,
    doi=true,
    eprint=false,
    hyperref=true,
    backref=false,
    parentracker=true,
    dashed=false,
    block=space
]{biblatex}
\addbibresource{referencias.bib}

% --- PAQUETES PARA DISEÑO, FORMATO Y CONTENIDO ---
\usepackage{geometry}
\geometry{a4paper, margin=2.54cm}
\usepackage{fancyhdr}
\usepackage{tabularx}
\usepackage{booktabs}
\usepackage{float}
\usepackage{caption}
\usepackage{listings}
\usepackage{enumitem}
\usepackage{amsmath}
\usepackage{xcolor}
\usepackage{setspace}
\usepackage{mathpazo} % Para la fuente Palatino

% --- PAQUETE HYPERREF (DEBE SER DE LOS ÚLTIMOS EN CARGARSE) ---
\usepackage{xurl} % Permite saltos de línea en URLs
\usepackage{hyperref}

% --- CONFIGURACIONES DE PAQUETES (DESPUÉS DE CARGARLOS) ---
\hypersetup{
    colorlinks=true,
    linkcolor=blue!50!black,
    citecolor=blue!50!black,
    urlcolor=cyan!50!black,
    pdftitle={Análisis de Procesos Organizacionales para un Sistema de Gemelo Digital de Mantenimiento},
    pdfauthor={Juan David Julio Serrano}
}

\newcolumntype{L}{>{\raggedright\arraybackslash}X} 

\definecolor{codegray}{rgb}{0.5,0.5,0.5}
\lstset{
    basicstyle=\ttfamily\small, breaklines=true, frame=single,
    showstringspaces=false, numbers=left, numberstyle=\tiny\color{codegray},
    keywordstyle=\color{blue},
    commentstyle=\color{green!40!black},
    stringstyle=\color{purple}
}

\captionsetup{
    format=plain,
    labelfont=bf,
    textfont=normalfont,
    singlelinecheck=false,
    justification=centering
}

% ===================================================================
%                       METADATOS DEL DOCUMENTO
% ===================================================================
\title{Análisis de Procesos Organizacionales para un Sistema de Gemelo Digital de Mantenimiento}
\author{Juan David Julio Serrano}
\date{12 de noviembre de 2025}

% ===================================================================
%                       INICIO DEL DOCUMENTO
% ===================================================================
\begin{document}

% --- PORTADA ---
\begin{titlepage}
    \begin{center}
        \vspace*{2.5cm}
        {\Huge \bfseries Identificación de Procesos Organizacionales para un Sistema de Gemelo Digital de Mantenimiento}\\[2cm]
        
        {\large \bfseries Aprendices:}\\[0.5cm]
        \large{Juan David Julio Serrano (Ficha: 3336237)}\\[0.3cm]
        \large{Maria Juliana Calvo Corredor (Ficha: 3336237)}\\[0.3cm]
        \large{Jesús David Algeciras Abril (Ficha: 3336238)}\\[0.3cm]
        \large{Didier Moreno Caicedo (Ficha: 3336238)}\\[1.2cm]
        
        {\large \bfseries Instructor:}\\[0.5cm]
        \large{Ing. Wilson Alberto Mosquera Caicedo}\\[1.5cm]
        
        \vspace{1cm}
        \large{Servicio Nacional de Aprendizaje (SENA)}\\[0.4cm]
        \large{Centro de Biotecnología Industrial - Regional Valle}\\[0.4cm]
        \large{Tecnología en Análisis y Desarrollo de Software (ADSO)}\\[1.2cm]
        \large{14 de noviembre de 2025}
    \end{center}
\end{titlepage}

% --- ÍNDICE DE CONTENIDOS ---
\clearpage
\tableofcontents
\clearpage
\doublespacing

% --- CONTENIDO DEL DOCUMENTO ---
\section{Introducción}

El presente documento tiene como objetivo aplicar los principios de la Teoría General de Sistemas (GST) para identificar y estructurar los elementos fundamentales de un \textbf{Sistema Computarizado de Gestión de Mantenimiento (CMMS)} o un sistema de \textbf{Gestión de Activos Empresariales (EAM)}, basado en un \textbf{Gemelo Digital Interactivo} para plantas industriales. Este ejercicio constituye la primera fase de análisis para el software a construir, cuyo propósito es optimizar la planificación, ejecución y análisis del mantenimiento de activos, centralizando la información y digitalizando los flujos de trabajo en campo.

Siguiendo la metodología de Análisis de Procesos de Negocio (BPA) propuesta en el material de estudio, la cual se fundamenta en la GST, este informe descompone el sistema en sus componentes clave: actores, procesos, entradas, salidas y relaciones \parencite{sena:fundamentos2024, ossa:2016}. Este informe busca establecer una base sólida y bien definida para las futuras fases de diseño y desarrollo, asegurando que la solución propuesta responda de manera coherente y completa a las necesidades de un entorno industrial moderno.

\section{Actores Responsables}

Los actores del sistema se dividen en tres categorías según su función dentro del ciclo de vida del mantenimiento:

\subsection{Roles Estratégicos y de Gestión}

\begin{itemize}
    \item \textbf{Gerente de Planta:} Usuario de alto nivel que consume información. Visualiza dashboards con Indicadores Clave de Rendimiento (KPIs) sobre la salud de los activos, costos y tiempos de inactividad para la toma de decisiones estratégicas \parencite{chen:2025, tractian:2025a, wireman:2013}.
    \item \textbf{Supervisor/Gestor de Mantenimiento:} Responsable de la estrategia general. Planifica los cronogramas de mantenimiento, gestiona presupuestos, aprueba, asigna órdenes de trabajo y supervisa al equipo operativo \parencite{accruent:2025, fracttal:2024, laurila:2017, tractian:2025b}.
    \item \textbf{Planificador de Mantenimiento:} Rol táctico que optimiza la programación de tareas, asegurando la disponibilidad de recursos como personal y repuestos para minimizar el impacto en la producción \parencite{fracttal:2024, laurila:2017, tractian:2025b}.
\end{itemize}

\subsection{Roles Operativos y de Campo}

\begin{itemize}
    \item \textbf{Técnico de Mantenimiento:} El usuario principal en campo. Utiliza una tablet para realizar inspecciones, ejecutar reparaciones, consultar documentación técnica y reportar el estado de las tareas.
    \item \textbf{Inspector HSEQ:} Especialista que utiliza la plataforma para realizar inspecciones de Activos Críticos de Seguridad (ACS), reportar anomalías y verificar su correcta resolución.
    \item \textbf{Jefe de Almacén:} Encargado de la gestión del inventario de repuestos y materiales de mantenimiento (conocido como inventario MRO). Procesa las solicitudes del sistema, actualiza el stock y registra la entrega de materiales.
    \item \textbf{Contratista Externo:} Usuario con acceso limitado que recibe órdenes de trabajo especializadas, ejecuta la labor y sube su informe para su procesamiento en el sistema \parencite{fracttal:2024, honeywell:2018, laurila:2017, servicenow:sf, tractian:2025b, wireman:2013}.
\end{itemize}

\subsection{Roles de Soporte Técnico y Analítico}

\begin{itemize}
    \item \textbf{Ingeniero de Proyectos (IT/OT/CAD):} Rol multifuncional responsable del Gemelo Digital. Se encarga de crear y mantener el modelo 3D (CAD), gestionar la infraestructura de software (IT) y asegurar la integración con sensores y maquinaria física (OT).
    \item \textbf{Analista de Datos y Fiabilidad:} Utiliza datos históricos y de sensores para desarrollar modelos predictivos, analizar tendencias de fallos, calcular métricas como el \textbf{Tiempo Medio Entre Fallos (MTBF)} y proponer mejoras en la estrategia de mantenimiento.
    \item \textbf{Administrador del Sistema:} Rol técnico que gestiona usuarios, permisos, integraciones con otros sistemas (ej. ERP) y realiza el mantenimiento general de la plataforma \parencite{chen:2025, ibm:2024, laurila:2017, mrzyk:2023, osullivan:2020, stum:2018, zhang:2025}.
\end{itemize}

\section{Procesos y Subprocesos}

De acuerdo con la estructura de una organización enfocada en el mantenimiento industrial, un sistema de software efectivo se articula en torno a macroprocesos que van desde la gestión de la información fundamental hasta la ejecución de estrategias de mantenimiento \parencite{ibm:2024, stum:2018}. Herramientas como los sistemas EAM o CMMS son fundamentales para centralizar y automatizar estas funciones, con el fin de optimizar el rendimiento y la vida útil de los activos \parencite{fluke:2020, sap:sf, tractian:2025a}.

\subsection{Procesos de Soporte y Gestión de Información}

\begin{enumerate}
    \item \textbf{Gestión de Activos:} Proceso fundamental que consiste en crear y mantener un registro centralizado de todos los activos de la planta, incluyendo sus componentes, manuales técnicos, historial y garantías. Este repositorio único es la piedra angular de cualquier sistema de mantenimiento moderno \parencite{ftmaintenance:2019, manwinwin:2015, mazenko:2015, servicenow:sf}.
    \item \textbf{Gestión de Inventario (MRO):} Controla el stock de piezas de repuesto, herramientas y otros materiales (Maintenance, Repair, and Operations). Este módulo incluye la gestión de compras y la asignación de recursos a las órdenes de trabajo, siendo un componente central en la mayoría de las plataformas CMMS/EAM \parencite{ftmaintenance:2019, honeywell:2018, laurila:2017, stum:2018}.
    \item \textbf{Visualización e Interacción con Gemelo Digital:} Permite a los usuarios navegar por una representación 2D/3D de la planta y activar capas de información (eléctrica, mecánica, HSEQ). Al seleccionar un activo, se accede a su "hoja de vida" digital, que es una representación dinámica que emula el rendimiento de su contraparte física y puede incluir dashboards de salud en tiempo real, historial de intervenciones y vistas explosionadas de sus componentes \parencite{chen:2025, miao:2025, osullivan:2020}.
    \item \textbf{Generación de Informes y Análisis (KPIs):} El sistema procesa de forma continua los datos operativos y de mantenimiento para generar dashboards e informes automatizados. Estos informes presentan métricas clave para la gestión, como el Tiempo Medio Entre Fallos (MTBF), el Tiempo Medio de Reparación (MTTR) y la Efectividad General del Equipo (OEE), permitiendo un análisis basado en datos para la toma de decisiones gerenciales \parencite{fiix:2019, ftmaintenance:2019, hxgn:sf, obrien:sf}.
\end{enumerate}

\subsection{Procesos de Ejecución de Mantenimiento}

\begin{enumerate}
    \item \textbf{Ciclo de Mantenimiento Preventivo:} Este proceso proactivo se inicia cuando el Planificador de Mantenimiento programa órdenes de trabajo (OT) recurrentes, basándose en intervalos de tiempo fijos (ej., "inspección trimestral") o en métricas de uso del equipo. Una de las capacidades centrales de un CMMS es generar y asignar automáticamente estas órdenes de trabajo según el calendario o los umbrales de uso establecidos, asegurando que el mantenimiento se realice antes de que ocurra una falla \parencite{accruent:2025, collabit:2020, fmx:2019, honeywell:2018, laurila:2017, mazenko:2015}.
    \item \textbf{Ciclo de Mantenimiento Correctivo (Reactivo):} Este ciclo se activa tras la detección de un fallo y comienza con una \textbf{solicitud de trabajo (SR)}, que es la notificación inicial de una necesidad. El Supervisor revisa, aprueba y prioriza esta solicitud, convirtiéndola en una \textbf{orden de trabajo (OT) formal}. Posteriormente, el Planificador asigna los recursos necesarios, incluyendo el técnico y los repuestos. Al finalizar la reparación, el Técnico documenta el cierre, registrando información crucial para el historial del activo, como los códigos de falla, la descripción de la solución, las horas-hombre empleadas y los materiales utilizados \parencite{accruent:2025, fracttal:2024, landau:2023, stum:2018, tractian:2025b}.
    \item \textbf{Ciclo de Mantenimiento Predictivo (PdM):} Esta estrategia avanzada se fundamenta en la monitorización en tiempo real del estado de los activos mediante sensores IoT (ej. vibración, temperatura). Los Gemelos Digitales (DT), potenciados con modelos de Machine Learning (ML), procesan estos datos para pronosticar fallos inminentes y estimar la vida útil restante (RUL) de los componentes. El objetivo es transformar la gestión de mantenimiento de un enfoque reactivo a uno \textbf{proactivo}. El sistema genera una \textbf{alerta predictiva} que desencadena automáticamente una orden de trabajo, permitiendo una intervención justo a tiempo, maximizando la vida útil del activo y minimizando el tiempo de inactividad no planificado \parencite{chen:2025, ibm:2024, osullivan:2020, tractian:2025a}.
\end{enumerate}

\subsection{Procesos de Interoperabilidad}

\begin{enumerate}
    \item \textbf{Procesamiento Automatizado de Informes (LLM):} Para optimizar la gestión de datos no estructurados, el sistema utiliza un Large Language Model (LLM). Cuando un Contratista Externo sube un informe de trabajo (PDF, imagen), el LLM extrae la información clave (piezas, horas, descripción) y la presenta en un formato estandarizado para la aprobación del Supervisor. Esta capacidad se alinea con el uso de la Inteligencia Artificial Generativa (GenAI) para automatizar tareas complejas y repetitivas en el mantenimiento moderno \parencite{chen:2025, hxgn:sf}.
    \item \textbf{Sincronización con Sistemas Externos:} Se contempla la integración del sistema con plataformas empresariales existentes para asegurar la consistencia de los datos en toda la organización. La conexión con sistemas como ERP (Enterprise Resource Planning) es fundamental para unificar los flujos de información entre mantenimiento, finanzas y la cadena de suministro. Esta capacidad, denominada \textbf{"Interacción del Sistema"}, se realiza a través de APIs y es clave para un ecosistema EAM verdaderamente integrado \parencite{anylogic:sf, ibm:2024, laurila:2017, tractian:2025a, zhang:2025}.
\end{enumerate}

\section{Entradas y Salidas}

El sistema opera mediante un flujo continuo de datos que, una vez procesados, producen información de valor para la toma de decisiones. Estos flujos son la base de un CMMS/EAM efectivo, donde la calidad de la entrada de datos es crucial para el éxito del sistema \parencite{fiix:2019, fluke:2020, mazenko:2015, obrien:sf}.

\subsection{Entradas (Inputs)}

Son los datos que alimentan el sistema desde diversas fuentes para construir una imagen completa del estado de los activos y las operaciones.

\begin{itemize}
    \item \textbf{Datos de Planificación y Configuración:}
    \begin{itemize}
        \item \textbf{Plan de Mantenimiento Programado:} Calendario de inspecciones y mantenimientos preventivos (PM), una función esencial de cualquier CMMS que se basa en el tiempo o en métricas de uso como lecturas de medidores \parencite{accruent:2025, ftmaintenance:2019, laurila:2017}.
        \item \textbf{Modelo 3D/CAD del Activo:} Archivos de diseño que sirven como base visual para la creación y actualización del Gemelo Digital \parencite{miao:2025, zhang:2025}.
    \end{itemize}
    \item \textbf{Datos Operativos y de Campo:}
    \begin{itemize}
        \item \textbf{Datos de Sensores en Tiempo Real:} Lecturas continuas de variables físicas (vibración, temperatura, presión) desde los activos, siendo la base del Mantenimiento Basado en Condición (CBM) y Predictivo (PdM) \parencite{chen:2025, mrzyk:2023, osullivan:2020, sap:sf}.
        \item \textbf{Datos de Operación del Activo:} Registros de uso, como horas de funcionamiento y ciclos de producción, utilizados a menudo para programar el mantenimiento \parencite{laurila:2017, manwinwin:2015, osullivan:2020}.
        \item \textbf{Datos de Inspección y Anomalías:} Información capturada por los técnicos en campo a través de formularios o checklists digitales, incluyendo evidencia multimedia (fotos, videos) que documentan fallos o desviaciones \parencite{fracttal:2024, ftmaintenance:2019, servicenow:sf}.
    \end{itemize}
    \item \textbf{Datos de Transacción:}
    \begin{itemize}
        \item \textbf{Solicitud de Repuesto:} Petición formal de un componente específico del inventario, esencial para la gestión de MRO \parencite{hxgn:sf, stum:2018}.
        \item \textbf{Informe de Contratista:} Documentos no estructurados (PDF, JPG) con el reporte de un trabajo realizado por un tercero, que sirven como entrada para su procesamiento y estandarización \parencite{hxgn:sf, zhang:2025}.
    \end{itemize}
\end{itemize}

\subsection{Salidas (Outputs)}

Son los resultados generados por el sistema tras procesar las entradas, convirtiendo datos brutos en información accionable.

\begin{itemize}
    \item \textbf{Resultados de Proceso y Tareas:}
    \begin{itemize}
        \item \textbf{Orden de Trabajo Digital (OT/WO):} Documento formal que es el corazón del mantenimiento, detallando la tarea a ejecutar, los recursos necesarios, el responsable asignado y las instrucciones. Sirve para formalizar, planificar y registrar todo el trabajo \parencite{accruent:2025, landau:2023, stum:2018, tractian:2025b}.
        \item \textbf{Informe de Cierre de Trabajo:} Registro final de una OT completada, que documenta crucialmente las horas-hombre, piezas usadas y la descripción de la solución para análisis futuros \parencite{ardila-marin:2018, fracttal:2024, laurila:2017}.
        \item \textbf{Historial de Activo Actualizado:} Registro cronológico de todas las intervenciones realizadas sobre un activo, conformado por las órdenes de trabajo cerradas. Es la base para el análisis de costos y rendimiento a largo plazo \parencite{ardila-marin:2018, fracttal:2024}.
        \item \textbf{Documento Estructurado:} Información clave extraída y formateada a partir del informe de un contratista mediante el procesamiento del LLM.
    \end{itemize}
    \item \textbf{Resultados de Análisis y Notificaciones:}
    \begin{itemize}
        \item \textbf{Alertas y Notificaciones:} El sistema genera alertas automáticas que pueden ser operativas (ej., "stock bajo de un repuesto") o predictivas (ej., "alta probabilidad de fallo inminente"). Estas notificaciones son clave para una gestión proactiva \parencite{chen:2025, hxgn:sf, mazenko:2015, osullivan:2020}.
        \item \textbf{Dashboard de Salud del Activo:} Visualización en tiempo real del estado operativo y la degradación de los componentes de un activo, a menudo representado como un "Índice de Salud" \parencite{chen:2025, mrzyk:2023}.
        \item \textbf{Informe de KPIs:} Reportes y dashboards que presentan métricas clave para la gestión estratégica, tales como el MTBF, el MTTR y la OEE \parencite{fiix:2019, ftmaintenance:2019, obrien:sf, osullivan:2020}.
    \end{itemize}
\end{itemize}


\section{Relaciones Principales del Sistema}

Las interacciones entre los diferentes elementos del sistema definen su comportamiento y flujo de trabajo. Estas relaciones se pueden clasificar en las siguientes categorías:

\subsection{Relaciones Jerárquicas y de Asignación}

\begin{itemize}
    \item Un \textbf{Supervisor de Mantenimiento} gestiona a uno o varios \textbf{Técnicos}.
    \item Un \textbf{Técnico de Mantenimiento} (o un \textbf{Contratista Externo}) es asignado a una o muchas \textbf{Órdenes de Trabajo}.
\end{itemize}

\subsection{Relaciones de Generación de Procesos (Causa-Efecto)}

\begin{itemize}
    \item Un \textbf{Reporte de Anomalía} (originado por una inspección) genera una \textbf{Orden de Trabajo} de tipo \textit{Correctivo}.
    \item Una \textbf{Alerta Predictiva} (originada por el modelo de IA) genera una \textbf{Orden de Trabajo} de tipo \textit{Predictivo}.
    \item Un \textbf{Plan de Mantenimiento Programado} genera \textbf{Órdenes de Trabajo} de tipo \textit{Preventivo}.
\end{itemize}

\subsection{Relaciones Funcionales y de Datos}

\begin{itemize}
    \item Un \textbf{Activo} posee un único \textbf{Historial de Mantenimiento}, el cual se compone de múltiples \textbf{Órdenes de Trabajo} cerradas.
    \item Un \textbf{Activo} está asociado a uno o varios flujos de \textbf{Datos de Sensores}.
    \item Una \textbf{Orden de Trabajo} consume uno o varios \textbf{Repuestos} del \textbf{Inventario (MRO)}.
    \item El \textbf{Gerente de Planta} consulta los \textbf{Informes de KPIs}, los cuales son calculados a partir de los datos agregados de todas las \textbf{Órdenes de Trabajo} completadas.
\end{itemize}

\section{Contexto y Fronteras del Sistema}

En línea con el enfoque de la Teoría General de Sistemas, es fundamental definir el contexto y las fronteras que delimitan el alcance del software a construir.

\begin{itemize}
    \item \textbf{Contexto:} El sistema opera dentro del contexto más amplio de la \textbf{gestión de activos y operaciones de una planta industrial}. Su propósito es interactuar con los procesos de mantenimiento para optimizar la fiabilidad y productividad de los equipos.
    \item \textbf{Fronteras:} Las fronteras del sistema se definen por sus puntos de interacción:
    \begin{itemize}
        \item \textbf{Humanas:} Los \textbf{Actores Responsables} (Técnicos, Supervisores, Gerentes, etc.) son las interfaces humanas que operan dentro del sistema o interactúan con él.
        \item \textbf{Técnicas:} Los \textbf{Procesos de Interoperabilidad} definen las fronteras con otros sistemas de software. El sistema se conectará con plataformas ERP/EAM existentes y con herramientas de productividad (Google Workspace/Office 365), pero no es responsable de las funcionalidades internas de dichos sistemas. Su límite está en el intercambio de datos con ellos.
    \end{itemize}
\end{itemize}



% --- BIBLIOGRAFÍA ---
\clearpage
\printbibliography[heading=bibintoc,title=Referencias]

\end{document}